% Ad Atticum 8.15 --- headnote
% ad-atticum-08-15, 3 March 49 BC, Formian villa

\textit{Cicero to Atticus, written from the Formian
villa on the fifth day before the Nones of March 49 BC
(the manuscript dateline: \textit{Scr.\ in Formiano v
Non.\ Mart.\ a.\ 705 (49)}). Three letters from Atticus
have arrived together by the freedman Aegypta: an old
one from the fourth day before the Kalends (delayed
because Pinarius the carrier never reached Cicero), and
two from the day before the Kalends. Section 1 walks
through the first letter point by point ---
Vibullius's mission (Caesar did not even see him),
Atticus's question about how Cicero will receive Caesar
on his return (Cicero is thinking of avoiding him
altogether), Atticus's own plans to leave Italy for
Athens, and the still-open question whether Domitius is
travelling with his fasces.}

\textit{Section 2 turns to the two later letters and to
the central question --- to stay or to cross --- and
gives one of the bleakest formulations of the dilemma
in the whole correspondence: \textit{cautior certe est
mansio, honestior existimatur traiectio}. Staying is
safer; crossing is judged more honourable. Cicero would
sooner have many think him incautious than a few think
him dishonourable. The legal point is technical and
characteristic: nearly everyone in the Pompeian camp
has the right to be out of Italy --- they hold
\textit{imperium} or are legates of those who do, and
by ancestral custom the consuls may enter every
province --- so the crossing is regular in form even
when it is desertion in fact. The letter closes with a
note that he has sent on a copy of Cornelius Balbus's
letter (the Caesarian agent in Rome) so that Atticus
can grieve with him over being made a laughingstock.}
